\section{Contoh}
\label{sec:contoh-bab10}

Dua contoh berikut memperlihatkan operasi aritmatika lengkap dan
pemeriksaan hasilnya. Contoh pertama adalah program utuh yang
menyimpan hasil ke segmen data. Contoh kedua memperlihatkan
perkalian dan pembagian 16 bit dengan pasangan \reg{DX}:\reg{AX}.

\subsection{Contoh 1: Operasi Aritmatika Lengkap}

Program pada Kode~\ref{lst:operasi-aritmatika} memakai \asm{ADD},
\asm{ADC}, \asm{SUB}, \asm{SBB}, \asm{INC}, \asm{DEC}, \asm{MUL},
dan \asm{DIV} dalam satu program.

\lstinputlisting[
    language=MASM,
    caption={Program \file{operasi-aritmatika.asm} untuk mengamati
    operasi aritmatika dan \eng{flag}.},
    label={lst:operasi-aritmatika}
]{\CodeDir/operasi-aritmatika.asm}

Program tidak mencetak keluaran. Hasil setiap operasi disimpan pada
segmen data dan diperiksa dengan DEBUG atau emu8086. Bagian penting
program tersebut adalah sebagai berikut.

\begin{itemize}
    \item \lstasm|ADD AX, b_lo| menjumlahkan word rendah. Karena
          hasilnya melewati 16 bit, \flag{CF} bernilai 1.
    \item \lstasm|ADC AX, b_hi| menjumlahkan word tinggi ditambah
          \flag{CF}, sehingga hasil 32 bit menjadi
          \hex{0004}:\hex{0000}.
    \item \lstasm|SUB AX, d_lo| menghasilkan pinjaman, lalu
          \lstasm|SBB AX, d_hi| melanjutkan pengurangan dengan
          pinjaman tersebut.
    \item \lstasm|MUL BL| mengalikan \reg{AL} dengan \reg{BL} dan
          menaruh hasil 16 bit pada \reg{AX}.
    \item \lstasm|CMP AH, BL| memeriksa apakah hasil bagi \asm{DIV}
          akan muat pada \reg{AL}. Bila tidak, program melompat ke
          label \code{luapan}.
\end{itemize}

Hasil yang diharapkan dirangkum pada Tabel~\ref{tab:hasil-program}.

\begin{table}[H]
    \centering
    \caption{Hasil operasi pada \file{operasi-aritmatika.asm}.}
    \label{tab:hasil-program}
    \small
    \begin{tabular}{@{}p{4.4cm}p{2.6cm}p{6.2cm}@{}}
        \toprule
        \textbf{Operasi} & \textbf{Hasil} &
        \textbf{Letak hasil} \\
        \midrule
        \hex{0001FFFF} + \hex{00020001} &
        \hex{0004}:\hex{0000} &
        \code{hasil\_hi}:\code{hasil\_lo} \\
        \addlinespace
        \hex{00050000} $-$ \hex{00020001} &
        \hex{0002}:\hex{FFFF} &
        \code{kurang\_hi}:\code{kurang\_lo} \\
        \addlinespace
        $25 \times 4$ & \hex{0064} & \code{kali} \\
        \addlinespace
        $1000 \div 7$ & \hex{008E}, sisa \hex{0006} &
        \code{bagi}, \code{sisa} \\
        \bottomrule
    \end{tabular}
\end{table}

\subsection{Contoh 2: MUL dan DIV 16 Bit}

Perkalian 16 bit menghasilkan 32 bit pada \reg{DX}:\reg{AX}.
Pembagian 16 bit membagi \reg{DX}:\reg{AX} dengan operand dan
menyimpan hasil bagi di \reg{AX} serta sisa di \reg{DX}.

\begin{lstlisting}[language=MASM, numbers=none]
; MUL 16 bit: 1000 x 20 = 20000 (4E20h)
MOV AX, 1000
MOV BX, 20
MUL BX              ; DX:AX = 0000:4E20h

; DIV 16 bit: 100000 / 1000 = 100
MOV DX, 1
MOV AX, 86A0h       ; DX:AX = 0001:86A0h = 100000
MOV BX, 1000
CMP DX, BX          ; periksa agar hasil bagi muat di AX
JAE luapan
DIV BX              ; AX = 100, DX = 0
luapan:
\end{lstlisting}

Pada pembagian 16 bit, syarat aman adalah \reg{DX} lebih kecil
daripada pembagi. Bila \reg{DX} sama atau lebih besar, hasil bagi
tidak muat pada \reg{AX} dan terjadi \eng{divide overflow}.

\begin{infobox}
\vspace{1.1em}
Nilai \hex{4E20} adalah 20.000 desimal. Angka ini masih muat pada
\reg{AX}, sehingga \reg{DX} bernilai nol. Bila hasil perkalian
melebihi 65.535, \reg{DX} tidak nol dan \flag{CF} = \flag{OF} = 1.
\end{infobox}
